% page 5 of the facsimile (image p-004 of the scan)
% block 154.9 x 203.1 mm ; left margin 47.1 mm
\pgc{25.400mm}{0.000mm}{
\l{\cel{20}P\cel{2}R\cel{2}E\cel{2}F\cel{2}A\cel{2}C\cel{2}O}
\l{}
\l{\cel{22}----------------}
\l{}
\l{}
\l{\cel{5}Dicar ke interkompreno povas eventar exkluzive segun}
\l{quante la konversanti - lore aktiva, lore pasiva - atribuas}
\l{a la vorti quin li uzas senco\cel{1}\sou{unika}\cel{1}e\cel{1}\sou{perfekte definita}}
\l{esas truismo.}
\l{}
\l{\cel{5}Che la lingui nacionala, interkompreno esas desfacila}
\l{pro ke, en multa kazi, granda quanto de lia vorti havante}
\l{plura senci, un ek la konversanti atribuas (pro oblivio di}
\l{ta plurasenceso, o pro trompo-volo) a ta o ca vorto diskuto-}
\l{precipua quan lu uzas senco qua ne koincidas kun la senco}
\l{expektata da la altra persono.}
\l{}
\l{\cel{5}Quale evitar lo ?\cel{2}La pensisti maxim eminenta, en}
\l{omna epoki, respondis, racionale : \textquotedbl{}per donar, a singla}
\l{vorto, senco qua esas unika\textquotedbl{}.}
\l{}
\l{\cel{5}Tala idealo ne esas realigebla che la lingui nacionala}
\l{(ecepte en la domeno di la cienci, ube la defini di la}
\l{termini eventas ante lia uzo - defino da la aktiva-uzero}
\l{dat-unesma) pro ke ta kolekturo de la vorti quan esas la}
\l{dicionario di ta o ca linguo indikas (por la personi qui uzas}
\l{ta linguo pasive) omna senci,ulakaze tre interdiferanta, quin}
\l{atribuas a ta o ca vorto la personi qui uzas olu aktive.}
\l{Pro lo, la problemo esas nesolvebla (adminime, lu ne esas}
\l{solvebla sen disipo di tempo, di energio e di pekunio).}
\l{}
\l{\cel{5}\sou{Kontraste}, la idealo esas realigita,\cel{1}\sou{fakte}, da la linguo}
\l{universala, supleanta e neutra, t.e.\cel{1}\sou{da Ido}\cel{1}- nam, che Ido,}
\l{}
\l{\cel{8}l. singla vorto havas senco\cel{1}\sou{unika}, e ta senco esas}
\l{\cel{11}definita precize, minucioze, - e}
\l{\cel{8}2. la personi qui uzas la linguo aktive agas lo}
\l{\cel{11}egardante (rigoroze ed exkluzive) la senco quan}
\l{\cel{11}donis a ta vorto la\cel{1}\sou{Akademio di Ido}, pleyado di}
\l{\cel{11}qua la membri esis maxime kompetenta pri linguis-}
\l{\cel{11}tiko, pri la cienci (kom exemplo: pri la psiko-}
\l{\cel{11}logio linguo-uzala) e pri la diversa faki di la}
\l{\cel{11}tekniki, e c.}
\l{}
\l{\cel{5}Ma ta realigo implikas ke, ultre lo jus dicita, omna}
\l{Ido-uzero, irgaloke en mondo, disponas\cel{1}\sou{kolekturo de ta de}-}
\l{\sou{fini},\cel{1}\sou{defini per Ido ipsa}.}
}
